% ===================================================================
%  도표 제 6 호 — 집 안, 세간, 먹을거리, 불 켜기.
%
%  Traduction, faite en 2026, en coreen d'aujourd'hui, sur l'ido de
%  texto/io. Regles de la colonne : voir 10-dopyo-01.tex. Cinq pieces,
%  cinq numerotations : les cles portent c1 a c5.
%
%  CINQUANTE-DEUX PAIRES, LE PLUS GRAND NOMBRE DU LIVRET, et la raison
%  est dans la nature meme du tableau : il ne raconte rien, il POSE
%  des objets les uns par rapport aux autres. Le seau DEVANT la table
%  de toilette, l'armoire QUI contient le linge, l'edredon SUR le
%  couvre-lit, le pot a fleurs DANS l'embrasure, le cadre PENDU au
%  mur, la marmite AVEC son couvercle, le tonneau MUNI de son
%  robinet. Chaque preposition est une inversion en coreen, et il y en
%  a cinquante-deux a retourner. Le tableau 5, deux fois plus long,
%  n'en avait que trente-sept : ce n'est decidement pas la longueur
%  qui coute.
%
%  LA MAISON COREENNE N'EST PAS CELLE DU LIVRET, ET LA COLONNE NE LA
%  DEPLACE PAS. Une maison coreenne de 1926 n'a ni lit a baldaquin, ni
%  cheminee a chenets, ni baignoire : on y dort sur le 온돌, on s'y
%  chauffe par le sol, on s'y lave au 목욕간. La tentation etait de
%  traduire 온돌 pour le lit et 아궁이 pour la cheminee — deux mots
%  coreens vivants et faux. La colonne ecrit 침대 le lit, 벽난로 la
%  cheminee, 욕조 la baignoire. Un seul 아궁이 est garde, au bloc
%  c4-01-1, parce que la LA chose est bien un foyer ouvert ou l'on
%  fait tourner une broche, et que le mot coreen le designe
%  exactement. On modernise la langue, jamais les choses — et cela
%  vaut aussi pour les choses qu'on aurait plaisir a retrouver.
%
%  CE QUE LE COREEN NOMME SANS AIDE, en revanche, remplit la cuisine
%  entiere : 부삽 la pelle a feu, 부집게 les pincettes, 장작받이 les
%  chenets, 재막이 le garde-cendres, 불가리개 l'ecran de foyer, 풀무
%  le soufflet, 석쇠 le gril, 도마 le billot, 강판 la rape, 깔때기
%  l'entonnoir, 개수대 l'evier, 국자 la louche, 거품 국자 l'ecumoire,
%  반짇고리 le panier a ouvrage, 골무 le de. Vingt-cinq objets de
%  cuisine, pas un emprunt.
%
%  LE RENVOI (6) EST DONNE DEUX FOIS PAR LE FAC-SIMILE — au savon puis
%  a la femme de chambre. C'est l'un des TROIS ECARTS DECLARES du
%  controle : le francais donne (16) a la femme de chambre et il a
%  raison, mais la source ido ne bouge pas, et la colonne porte ce que
%  la source porte.
% ===================================================================

%%K t06-apar-1 apar
\VUsaut{6.0mm}
\VUtitre{15.0pt}{60}{도표 제 6 호}
\VUsaut{1.4mm}
\VUtitre{11.4pt}{0}{집 안, 세간, 먹을거리,}
\VUsaut{0.4mm}
\VUtitre{11.4pt}{0}{불 켜기.}
\VUsaut{2.2mm}

%%K t06-apar-2 apar
집이 다 되었다. 요안네스의 아저씨는 새 집에 들었고, 우리는 그 짜임을 이미 안다. 이제 그가
집을 어떻게 꾸몄는지 보자. 먼저 둘러볼 곳은 이러하다.

%%K t06-tit-1 sub
\VUpk{10.2pt}{40}{침실.}
\VUsaut{3.0mm}

%%K t06-c1-01-1 p
1. --- 우리는 때를 잘못 맞추어 왔다. 하녀가 \VUgras{방}을 치우느라 바쁘기 때문이다.

%%K t06-c1-02-1 p
2. --- \VUgras{세면대} \textsuperscript{(1)} 위에는 씻고 빗고, 한마디로 몸단장하는 데
쓰는 여러 가지가 여기저기 놓여 있다. \VUgras{대야} \textsuperscript{(2)}와
\VUgras{물주전자} \textsuperscript{(3)}, \VUgras{머리솔} \textsuperscript{(4)}, \VUgras{빗}
\textsuperscript{(5)}, \VUgras{비누} \textsuperscript{(6)}와 \VUgras{스펀지}
\textsuperscript{(7)}, 끝으로 물치약을 담은 \VUgras{작은 병} \textsuperscript{(8)}이다.

%%K t06-c1-03-1 p
3. --- 세면대 앞 바닥에는 더러운 물을 받는 \VUgras{물통} \textsuperscript{(9)}이 있다.

%%K t06-c1-04-1 p
4. --- \VUgras{앞방} \textsuperscript{(10)}에는 \VUgras{옷걸이} \textsuperscript{(13)} 몇 개와
\VUgras{우산} \textsuperscript{(11)} 두 개가 보인다 — \VUgras{우산꽂이} \textsuperscript{(12)}
안에 꽂혀 있다.

%%K t06-c1-05-1 p
5. --- 문 반대쪽에는 \VUgras{거울장} \textsuperscript{(14)}이 있다. 그 안에
\VUgras{속옷가지} \textsuperscript{(15)}가 들어 있다.

%%K t06-c1-06-1 p
6. --- \VUgras{하녀} \textsuperscript{(6)}가 \VUgras{침대} \textsuperscript{(17)}를 매만지는
동안, 그 여러 부분을 살펴보자. \VUgras{침대틀} \textsuperscript{(20)}에는 좋은
\VUgras{용수철 요} \textsuperscript{(21)}가 들어 있고, 그 위에 유스티나가 곧 털
\VUgras{요} \textsuperscript{(22)}를 깔 것이다. 그 다음에 걸상에서 \VUgras{홑이불}
\textsuperscript{(23)} 둘과 \VUgras{이불} \textsuperscript{(24)}을 집을 것이다. 그러고 나서
머리맡에 \VUgras{긴 베개} \textsuperscript{(25)}와 \VUgras{베개} \textsuperscript{(26)}를
놓을 것이다. 그 베개들은 \VUgras{오리털 이불} \textsuperscript{(27)}과 함께
\VUgras{침대 깔개} \textsuperscript{(32)} 위에 있다. 그 여자는 먼지떨이로 \VUgras{휘장}
\textsuperscript{(19)}과 \VUgras{침대 닫집} \textsuperscript{(18)}의 먼지를 턴다. 이렇게
일의 한 몫이 끝난다.

%%K t06-c1-07-1 p
7. --- 그 다음에는 \VUgras{머리맡 탁자} \textsuperscript{(28)}를 좀 정돈하고,
\VUgras{자명종} \textsuperscript{(29)}을 다시 감고, \VUgras{밤등} \textsuperscript{(30)}을
채비하고, 촛대를 닦으려고 \VUgras{초} \textsuperscript{(31)}를 빼야 한다.

%%K t06-tit-2 sub
\VUsaut{5.0mm}
\VUpk{10.2pt}{40}{반짇고리와 벽난로 딸림것.}
\VUsaut{3.0mm}

%%K t06-c1-08-1 p
8. --- \VUgras{반짇고리} \textsuperscript{(34)}도 정돈이 몹시 필요하다 — \VUgras{발판}
\textsuperscript{(33)} 곁에 있는 것이다. 그 여자는 \VUgras{수놓은 것}
\textsuperscript{(35)}을 개어 \VUgras{반짇상} \textsuperscript{(38)} 안에 넣어야 한다.
\VUgras{골무}, \VUgras{실} \textsuperscript{(36)}, \VUgras{바늘} \textsuperscript{(37)},
\VUgras{가위} \textsuperscript{(39)}도 함께. 그러고는 \VUgras{열쇠} \textsuperscript{(40)}로
상 뚜껑을 잠가야 한다.

%%K t06-c1-09-1 p
9. --- 방 한가운데 \VUgras{외다리 탁자} \textsuperscript{(41)} 위에서, 유스티나는
\VUgras{물병} \textsuperscript{(42)}의 물을 갈 것이다. \VUgras{설탕}을 \VUgras{설탕 그릇}
\textsuperscript{(43)}에 담고, \VUgras{설탕 집게} \textsuperscript{(44)}를 닦고,
\VUgras{옷솔} \textsuperscript{(46)}을 치울 것이다.

%%K t06-c1-10-1 p
10. --- 방 오른쪽은 그 여자에게 일을 덜 시킬 것이다. \VUgras{긴 안락의자}
\textsuperscript{(47)}와 그 \VUgras{방석} \textsuperscript{(48)}이 제자리에 있고, 추운
철이 아니어서 \VUgras{벽난로} \textsuperscript{(49)}의 딸림것들도 하나도 흐트러지지 않았기
때문이다. \VUgras{부삽} \textsuperscript{(50)}, \VUgras{부집게} \textsuperscript{(51)},
\VUgras{장작받이} \textsuperscript{(55)}, \VUgras{재막이} \textsuperscript{(56)}가 깨끗하다.
\VUgras{불가리개} \textsuperscript{(54)}도 옮겨지지 않았고, \VUgras{장작 궤}
\textsuperscript{(52)}에는 \VUgras{장작} \textsuperscript{(53)}이 넉넉히 채워져 있다.

%%K t06-noto-1 noto
\VUnotes{143.2mm}{%
(*) Babo = 어린아이; 영어 \textit{baby}, 프랑스어 \textit{bébé}, 이탈리아어
\textit{bambino}.%
}

%%K t06-noto-2 noto
\VUnotes{143.2mm}{%
(*) Lambrequino = 프랑스어 \textit{lambrequin}, 스페인어 \textit{lambrequines},
포르투갈어 \textit{lambrequins}, 독일어와 영어에서도 \textit{lambrequins};
그러므로 독·영·프·포·스.%
}

%%K t06-c1-11-1 p
11. --- 그래도 그 여자는 \VUgras{추시계} \textsuperscript{(57)}와 \VUgras{촛대}
\textsuperscript{(58)}와 \VUgras{잔물건 그릇} \textsuperscript{(59)}의 먼지를 털고, 끝으로
\VUgras{거울} \textsuperscript{(60)}을 닦아야 한다.

%%K t06-c1-12-1 p
12. --- 어린아이(아기 (*))의 \VUgras{요람} \textsuperscript{(61)}은 이미 채비가 되어 있다.

%%K t06-tit-3 sub
\VUsaut{5.0mm}
\VUpk{10.2pt}{40}{응접실.}
\VUsaut{3.0mm}

%%K t06-c2-01-1 p
1. --- 이제 응접실이다. 아주 고운 \VUgras{방}이다. 오른쪽 구석에 있는 벽난로는 섬세한
\VUgras{작은 조각상} \textsuperscript{(1)}과 고운 \VUgras{꽃병} \textsuperscript{(3)} 둘로
꾸며져 있고, 벨벳 \VUgras{휘장 장식} \textsuperscript{(2)} (*)이 드리워져 있다.

%%K t06-c2-02-1 p
2. --- 화덕의 불티가 깔개를 태우지 않도록, 화덕 앞에 구리 \VUgras{불티 가리개}
\textsuperscript{(4)}를 놓았다.

%%K t06-c2-03-1 p
3. --- 그 가까이에 맵시 있는 부인용 \VUgras{책상} \textsuperscript{(5)}이 열려 있다. 그
위에서 \VUgras{안주인} \textsuperscript{(23)}이 편지를 쓴다.

%%K t06-c2-04-1 p
4. --- 창은 \VUgras{창 커튼} \textsuperscript{(6)}으로 꾸며져 있고, \VUgras{매듭끈}
\textsuperscript{(8)}이 달려 \VUgras{커튼 고리} \textsuperscript{(7)}에 걸려 있다. 또
\VUgras{위에 드리운} \textsuperscript{(9)} 두꺼운 천 커튼도 있다.

%%K t06-c2-05-1 p
5. --- 창 안쪽 오목한 자리에는 \VUgras{화분 받이} \textsuperscript{(11)}가 있고, 그
안에 푸른 \VUgras{화초} \textsuperscript{(10)}가 담겨 있다. 훌륭한 종려나무인데, 그 잎이
\VUgras{석유등} \textsuperscript{(12)}에 닿을 만큼 뻗어 있다.

%%K t06-c2-06-1 p
6. --- 등의 \VUgras{등갓} \textsuperscript{(13)}은 마리아가 매만졌다. 마리아는 고운
\VUgras{젊은 여자} \textsuperscript{(14)}로, \VUgras{피아노} \textsuperscript{(15)} 앞에 있다.
\VUgras{걸상} \textsuperscript{(16)} 위에 아주 바르게 앉아 악보를 익힌다. 그 여자는 아주
솜씨 있고 꼼꼼한데, 완벽하게 정돈된 제 \VUgras{악보장} \textsuperscript{(17)}이 그것을
말해 준다.

%%K t06-tit-4 sub
\VUsaut{5.0mm}
\VUpk{10.2pt}{40}{개구쟁이.}
\VUsaut{3.0mm}

%%K t06-c2-07-1 p
7. --- 그 여자의 오라비 파울루스는 \VUgras{책} \textsuperscript{(19)}을 찾는다 —
\VUgras{책장} \textsuperscript{(18)}에서. 그는 \VUgras{젊은이} \textsuperscript{(20)}인데,
가만있지 못하고 견디기 어렵다. 너무 위에 있는 책에 닿으려고 책장에 매달리다가, 그 위에
있는 \VUgras{흉상} \textsuperscript{(21)}을 떨어뜨릴 뻔한다.

%%K t06-c2-08-1 p
8. --- 그가 그런 어리석은 짓을 하는 것은, 어머니가 벗과 이야기하느라 바쁜 틈을 탄
것이다. 그 벗은 며칠 묵으러 온 \VUgras{부인} \textsuperscript{(22)}이다. 어머니는
\VUgras{긴 의자} \textsuperscript{(24)}에 앉아 있는데, 그것은 \VUgras{안락의자}
\textsuperscript{(25)} 곁에 있다. 그 벗은 \VUgras{의자} \textsuperscript{(27)} 위에 있는데,
그 의자는 \VUgras{등받이} \textsuperscript{(26)}만 보인다.

%%K t06-c2-09-1 p
9. --- 벽에 걸린 큰 \VUgras{액자} \textsuperscript{(30)}에는 친척 아주머니의
\VUgras{초상} \textsuperscript{(29)}이 들어 있다. 상 위에 놓인 \VUgras{사진첩}
\textsuperscript{(32)}에는 나머지 식구들의 \VUgras{사진} \textsuperscript{(33)}이 모여 있다.
아이들이 방금 그것을 넘겨 보았다. \VUgras{의자} \textsuperscript{(31)}들이 아직 상 둘레에
모여 있기 때문이다.

%%K t06-c2-10-1 p
10. --- 사기로 된 \VUgras{중국 꽃병} \textsuperscript{(34)}은 \VUgras{종이칼}
\textsuperscript{(35)} 곁에 있는데, 마리아가 제 축일에 선물로 받은 것이다. 방 구석을
꾸미는 \VUgras{병풍} \textsuperscript{(36)}은 그 여자의 아저씨가 중국에서 가져온 것이다.

%%K t06-c2-11-1 p
11. --- 훌륭한 \VUgras{그림} \textsuperscript{(37)}이 걸려 있어 밝고 — \VUgras{칸막이벽}
\textsuperscript{(42)}에 걸려 있다 —, \VUgras{천장 등} \textsuperscript{(38)}이 비추어 주니
— 그 \VUgras{전등} \textsuperscript{(39)}이 저녁마다 즐겁게 빛난다 — 이 응접실은 아주
기분 좋아 보인다. 마루 위에 펼친 부드러운 \VUgras{깔개} \textsuperscript{(40)}와, 쓰기
그토록 쉬운 \VUgras{전기 초인종} \textsuperscript{(41)}이 그 아늑함을 마무리한다.

%%K t06-tit-5 sub
\VUsaut{3.6mm}
\VUpk{10.2pt}{40}{식당.}
\VUsaut{3.0mm}

%%K t06-c3-01-1 p
1. --- 저녁 종이 울렸다. 마리아만 빼고 온 식구가 상 둘레에 모였다. 식당은 사기로 된
훌륭한 \VUgras{난로} \textsuperscript{(1)} 덕분에 기분 좋게 데워져 있다. 가느다란
\VUgras{연통} \textsuperscript{(2)}이 달린 난로다.

%%K t06-c3-02-1 p
2. --- 이 방에는 세간이 많지 않다. 벽을 따라 \VUgras{걸상} \textsuperscript{(4)} 위에
꾸밈용 \VUgras{작은 수반} \textsuperscript{(3)}이 걸려 있다. 아주 가까이에 \VUgras{찬장}
\textsuperscript{(5)}이 있어, 그 위에 찬 먹을거리와 후식 접시와 곧 쓰지 않을 여러
상차림을 올려놓는다.

%%K t06-c3-03-1 p
3. --- 그리하여 윗칸에는 \VUgras{구운 닭} \textsuperscript{(7)}, \VUgras{생선}
\textsuperscript{(8)}, 그리고 \VUgras{굽접시} \textsuperscript{(10)}가 보인다 —
\VUgras{과일} \textsuperscript{(9)}이 담긴 것이다. 아래에는 \VUgras{치즈}
\textsuperscript{(11)}, \VUgras{빵} \textsuperscript{(12)}, \VUgras{양념병 받침}
\textsuperscript{(13)}이 있다. 그 받침에는 샐러드에 칠 것이 다 있다. 기름, 식초,
\VUgras{소금} \textsuperscript{(14)}, \VUgras{후추} \textsuperscript{(15)}다. 끝으로 아랫칸에
\VUgras{케이크} \textsuperscript{(17)}가 있다 — \VUgras{겨자 그릇} \textsuperscript{(16)}
곁에.

%%K t06-c3-04-1 p
4. --- 식당 시계는 \VUgras{벽시계} \textsuperscript{(20)}라고 하는데, 아주 남다른 꼴을
하고 있다. 나무 테에 끼워져 있고, 그 테에는 \VUgras{기압계} \textsuperscript{(19)}와
\VUgras{온도계} \textsuperscript{(18)}도 함께 들어 있다.

%%K t06-tit-6 sub
\VUsaut{4.6mm}
\VUpk{10.2pt}{40}{저녁 시중.}
\VUsaut{3.0mm}

%%K t06-c3-05-1 p
5. --- 방의 모든 벽에는 \VUgras{나무 널판} \textsuperscript{(21)}이 대어져 있다. 밀랍을
먹인 참나무다. 천으로 된 \VUgras{문휘장} \textsuperscript{(22)}이
베란다로 나가는 문을 넉넉히 가린다. 식구들은 저녁을 마치고 거기로 커피를 마시러 갈
것이다. \VUgras{찻잔} \textsuperscript{(24)}과 \VUgras{받침 접시} \textsuperscript{(25)}는
벌써 \VUgras{그릇장} \textsuperscript{(23)} 위에 채비되어 있다.

%%K t06-c3-06-1 p
6. --- 상 위에는 천장에 매단 \VUgras{매단 등} \textsuperscript{(26)}이 있는데, 그 아주
밝은 등이 저녁마다 식당 전체를 비춘다.

%%K t06-c3-07-1 p
7. --- 이제 \VUgras{밥상} \textsuperscript{(27)} 자체를 살펴보자. 식구들이 들어왔을 때
상차림은 다 되어 있었다. \VUgras{상보} \textsuperscript{(28)} 위, 함께 먹는 이마다 제
자리에 \VUgras{접시} \textsuperscript{(33)}, \VUgras{냅킨} \textsuperscript{(29)},
\VUgras{숟가락} \textsuperscript{(34)}, \VUgras{포크} \textsuperscript{(35)},
\VUgras{나이프} \textsuperscript{(36)}, \VUgras{잔} \textsuperscript{(32)}이 놓여 있었다.
포도주를 담은 \VUgras{병} \textsuperscript{(30)}과 물을 담은 \VUgras{물병}
\textsuperscript{(31)}도 있었다.

%%K t06-c3-08-1 p
8. --- \VUgras{하인} \textsuperscript{(39)}이 먼저 \VUgras{국그릇} \textsuperscript{(37)}을
날라 왔다. 그 안에는 아직 국이 김을 내고 있다. 지금 그는 첫 \VUgras{요리}
\textsuperscript{(38)}, 곧 고기 요리를 낸다. 다음에는 우리가 이미 이름을 댄 것들을
낼 것이고, 끝으로 \VUgras{후식}이 끝나면 주인들이 베란다로 들어간 틈을 타 상 도구를
치우고 식당을 다시 정돈할 것이다.

%%K t06-c3-08-2 p
\textit{알림.} --- \textit{먹을거리에 관한 흔한 낱말들은 여기에 아주 간추려서만
적어 두었다. 그 목록은 「호텔」(제 13 호)과 「시장」(제 15 호) 두 도표에서 채워질
것이다.}

%%K t06-tit-7 sub
\VUsaut{4.6mm}
\VUpk{10.2pt}{40}{부엌.}
\VUsaut{3.0mm}

%%K t06-c4-01-1 p
1. --- 반쯤 열린 문으로 들여다보는 부엌에는 그림 같은 어수선함이 있다. \VUgras{아궁이}
\textsuperscript{(1)} 앞에 — 그 안에서 \VUgras{불} \textsuperscript{(3)}이 타닥거린다 —
\VUgras{꼬치 돌리개} \textsuperscript{(5)}가 놓여 있다. 고운 \VUgras{구운 고기}
\textsuperscript{(7)}가 \VUgras{꼬치에 꿰여} \textsuperscript{(6)} 익어 가고 있다.

%%K t06-c4-02-1 p
2. --- 방금 쓴 \VUgras{풀무} \textsuperscript{(8)}와 \VUgras{숯삽} \textsuperscript{(9)}은
\VUgras{장작} \textsuperscript{(10)}과 마찬가지로 바닥에 놓인 채다.

%%K t06-c4-03-1 p
3. --- 벽난로 위는 정돈되어 있다. \VUgras{초롱} \textsuperscript{(11)},
\VUgras{성냥통} \textsuperscript{(13)} — 그 안에 \VUgras{성냥} \textsuperscript{(12)}이
있다 —, \VUgras{촛대} \textsuperscript{(14)}, \VUgras{촛대 받침} \textsuperscript{(15)}과
거기 꽂힌 \VUgras{초} \textsuperscript{(16)}가 제자리에 있다. 그러나 큰 \VUgras{숯통}
\textsuperscript{(18)}은 \VUgras{석탄} \textsuperscript{(17)}으로 가득한데, 그것을
\VUgras{찬모} \textsuperscript{(23)}가 방금 지하실에서 채워 온 뒤 부엌 한가운데에 그대로
두었다.

%%K t06-c4-04-1 p
4. --- 사실 그 딱한 아낙은 아침상이 제때 되도록 서둘러야 한다. 지금 그 여자는
\VUgras{부엌 화덕} \textsuperscript{(19)} 곁에서 \VUgras{소스} \textsuperscript{(24)}를
젓고 있다. 너무 태우면 안 되는 것이다.

%%K t06-c4-05-1 p
5. --- \VUgras{오븐} \textsuperscript{(20)} 안에서는 케이크가 구워진다. 아이들 새참으로
그 여자가 만든 것이다. 부엌 화덕 위에는 \VUgras{국자} \textsuperscript{(21)}와
\VUgras{거품 국자} \textsuperscript{(22)}가 걸려 있다.

%%K t06-tit-8 sub
\VUsaut{4.0mm}
\VUpk{10.2pt}{40}{그릇붙이.}
\VUsaut{3.0mm}

%%K t06-c4-06-1 p
6. --- 방 안쪽에 걸린 \VUgras{부엌 그릇} \textsuperscript{(25)}은 아주 깨끗하다. 고운
구리 \VUgras{냄비} \textsuperscript{(26)}, \VUgras{우유 단지} \textsuperscript{(27)},
\VUgras{작은 체} \textsuperscript{(28)}, \VUgras{강판} \textsuperscript{(29)}이 정성껏
닦여 있었다.

%%K t06-c4-07-1 p
7. --- \VUgras{소금 통} \textsuperscript{(30)}은 그릇장 위, \VUgras{차 주전자}
\textsuperscript{(31)}와 \VUgras{기름 큰 병} \textsuperscript{(32)} 곁에 놓여 있다.
\VUgras{서랍} \textsuperscript{(33)}에는 아마 상 도구와 부엌칼이 들어 있을 것이다.

%%K t06-c4-08-1 p
8. --- \VUgras{빗자루} \textsuperscript{(34)}와 \VUgras{먼지떨이} \textsuperscript{(35)}는
구석에 있다. 찬모가 방금 \VUgras{통나무 도마} \textsuperscript{(36)}를 썼는데, 그것을 창
곁에 두었다.

%%K t06-c4-09-1 p
9. --- \VUgras{손행주} \textsuperscript{(37)}가 벽에 걸려 있다 — \VUgras{깔때기}
\textsuperscript{(38)} 곁에. 부엌의 그쪽에는 \VUgras{석쇠} \textsuperscript{(39)},
\VUgras{식칼} \textsuperscript{(40)}, \VUgras{도마} \textsuperscript{(41)}가 치워져 있다.
그 다음, 식칼 위 \VUgras{선반} \textsuperscript{(42)}에는 \VUgras{단지}
\textsuperscript{(43)}와 \VUgras{냄비} \textsuperscript{(44)}가 있는데, 그 냄비에는
\VUgras{뚜껑} \textsuperscript{(45)}이 덮여 있다. \VUgras{큰솥} \textsuperscript{(46)}도
있다. \VUgras{프라이팬} \textsuperscript{(47)}은 그 가까이에 걸려 있다.

%%K t06-c4-10-1 p
10. --- 이 구석에서 빛을 내는 것은 구리 \VUgras{수도꼭지} \textsuperscript{(48)}뿐이다.
\VUgras{개수대} \textsuperscript{(49)}는 그 아래에 있는데, 그 위에 두꺼운
\VUgras{큰 물동이} \textsuperscript{(50)}가 놓여 있다. 개수대에는 작은 장이 딸려 있어 아주
쓸모 있어 보인다. 그 안에는 \VUgras{식초 통} \textsuperscript{(51)}을 두는데, 통에는
제 \VUgras{꼭지} \textsuperscript{(52)}가 달려 있다. \VUgras{찌꺼기 통}
\textsuperscript{(53)}과 \VUgras{더러운 그릇} \textsuperscript{(54)}도 거기에 있다.

%%K t06-tit-9 sub
\VUsaut{4.0mm}
\VUpk{10.2pt}{40}{부엌에 놓인 다른 것들.}
\VUsaut{3.0mm}

%%K t06-c4-11-1 p
11. --- \VUgras{큰 함지} \textsuperscript{(55)}는 \VUgras{남새} \textsuperscript{(58)}를
씻는 데 쓰는 것이고, 무쇠 \VUgras{솥} \textsuperscript{(56)}은 \VUgras{타일 바닥}
\textsuperscript{(57)} 위에 놓여 있다. 그 둘도 아마 그 장에 제자리가 있을 것이다.

%%K t06-c4-12-1 p
12. --- 찬모에게 화덕이 모자라지는 않는다. 상 모퉁이에 \VUgras{가스 화덕}
\textsuperscript{(59)}이 또 있어, 그 위 작은 냄비에서 물이 데워지고 있기 때문이다.
거기에서 새어 나오는 \VUgras{김} \textsuperscript{(60)}이 아마 물이 끓는다는 것을 알려
준다. 이 물은 아마 커피에 쓰일 것이다. 상 저쪽 끝에 \VUgras{커피 주전자}
\textsuperscript{(64)}와 \VUgras{커피 분쇄기} \textsuperscript{(63)}가 있으니 말이다.

%%K t06-c4-13-1 p
13. --- 화덕은 \VUgras{가스 꼭지} \textsuperscript{(62)}에 이어져 있다 — 긴
\VUgras{고무 관} \textsuperscript{(61)}으로.

%%K t06-c4-14-1 p
14. --- 찬모를 도우러 오는 \VUgras{날품 하녀} \textsuperscript{(66)}는 저녁상에 쓸 남새를
다듬는다. 다듬고 씻고 나면 그 여자는 그것을 \VUgras{함지} \textsuperscript{(65)}에 담아
삶을 때를 기다릴 것이다.

%%K t06-tit-10 sub
\VUsaut{4.0mm}
{\centering\fontsize{10.2pt}{10.2pt}\selectfont\textls[40]{\scshape 목욕실.} \textit{(목욕하는 방.)}\par}
\VUsaut{3.0mm}

%%K t06-c5-01-1 p
1. --- 이 집의 주요한 방들을 다 알려면 이제 한 가지 실례만 남았다. 목욕실의 꾸밈을 보는
일이다. 오래 걸리지는 않을 것이다. 넓지 않기 때문이다. \VUgras{욕조}
\textsuperscript{(1)}에 좋은 \VUgras{목욕 물끓이개} \textsuperscript{(2)}가 달려 있다는
것, \VUgras{비누 접시} \textsuperscript{(3)}가 목욕하는 이의 손에 닿는다는 것,
\VUgras{수건걸이} \textsuperscript{(5)}가 \VUgras{세면 수건} \textsuperscript{(4)}을 쉽게
말려 주리라는 것을 우리는 곧 알게 된다.

%%K t06-c5-02-1 p
2. --- 이 방은 벽이 \VUgras{타일} \textsuperscript{(6)}로 덮여 있다. 그 덕분에
\VUgras{샤워 장치} \textsuperscript{(7)}를 놓을 수 있었다. 쓸 때 튀는 물이 벽을 상하게
할까 걱정할 것 없이 말이다. 발 씻는 데 쓰는 아연 \VUgras{작은 대야} \textsuperscript{(8)}
가 세간을 마무리한다.

\VUsaut{6.0mm}
\VUfilet{22mm}
